\documentclass[11pt,a4paper]{article}
\usepackage[utf8]{inputenc}
\usepackage{mathtools, amssymb}
\usepackage{graphicx}
\graphicspath{{figures/}{data/drafts/figures/}{../}{data/}}
\usepackage{booktabs}
\usepackage{natbib}
\usepackage[margin=1in]{geometry}
\usepackage{physics}
\usepackage{xcolor}
\usepackage{authblk}
\usepackage{siunitx}
\usepackage{hyperref}

\title{Metric-Consistent Spectral Finite Element Decomposition of Stable Boundary Layer Regimes: \\
A Riemannian Manifold Approach to Sub-Mesoscale Wave--Turbulence Partitioning}

\author{David E. England\textsuperscript{1,*}\thanks{ORCID: \href{https://orcid.org/0009-0001-2095-6646}{0009-0001-2095-6646}}}
\affil{\textsuperscript{1}The University of Alabama in Huntsville, 301 Sparkman Drive, Huntsville, AL 35899, USA}
\affil{\textsuperscript{*}Corresponding author: David.England@UAH.Edu}

\date{\today}

\begin{document}

\maketitle

\begin{abstract}
Nocturnal stable boundary layers (SBLs) exhibit complex interactions between three dominant scales: large-scale mesoscale pressure forcing, horizontally propagating internal gravity waves (IGWs), and highly localized turbulent shear bursts. Traditional diagnostic frameworks conflate these regimes, producing high variance in inversion-layer flux estimates and inconsistent model-observation comparisons. We present a \emph{metric-consistent Riemannian pseudospectral} framework that projects sparse tower observations onto a non-uniform computational grid designed explicitly to maximize sensitivity near the surface (\SI{2.85}{\milli\meter} effective node resolution at $z = \SI{1.5}{\meter}$) while maintaining numerical stability at upper boundaries ($\kappa(\mathbf{M}) = 4619$). Using Chebyshev polynomials as the spectral basis and solving the rank-deficient least-squares problem via thin singular value decomposition (SVD), we extract three physically meaningful diagnostic features: effective modal dimension $D_{\mathrm{eff}}$, wave energy fraction $F_W$, and spectral curvature $\chi_N$.

Applying this framework to the CASES-99 campaign stable window (22--31 October 1999), we identify \emph{three statistically separable regimes} with a silhouette clustering score of $\bar{S} = 0.582$ for $K=3$: (1) Continuous Turbulence ($D_{\mathrm{eff}} > 20$, $F_W < 0.15$), (2) Wave-Dominated ($D_{\mathrm{eff}} \sim 4$--$6$, $F_W > 0.75$), and (3) Intermittent Shear Bursts ($2 < D_{\mathrm{eff}} < 15$, $0.15 < F_W < 0.50$). Crucially, the regime-specific rank compression ($\Delta_{\mathrm{rank}} = 24$) proves that atmospheric dynamics naturally occupy lower-dimensional manifolds and are not mathematical artifacts of basis truncation. 

We extend the theoretical framework by establishing an explicit physical link to the Ozmidov buoyancy scale $L_O$, resolving the grid-topology mismatch via a generalized metric-consistent Virtual Tower interpolation operator ($\mathcal{P}_{\mathrm{VT}}$) for direct NWP model validation, and formulating a $2/3$-rule spectral anti-aliasing filter for forward-predictive implementations. A synthetic gravity-wave reflection test confirms that the spectral partition functions suppress upper-boundary reflection by approximately $3.3\times 10^4$, validating the sponge-layer boundary conditions.
\end{abstract}

\noindent \textbf{Keywords:} stable boundary layers, spectral methods, manifold learning, gravity waves, Ozmidov scale, virtual tower interpolation

\newpage
\tableofcontents
\newpage

% ============================================================================
\section{Introduction}
\label{sec:intro}

Nocturnal stable boundary layers (SBLs) remain among the most poorly understood and inconsistently modeled features of the lower atmosphere \citep{Mahrt2014, Acevedo2016}. Unlike daytime convective boundary layers, where surface heating drives turbulence vertically, nocturnal SBLs exhibit a fundamental instability: strong temperature inversions suppress vertical mixing, yet persistent wind shear continues to generate turbulent kinetic energy (TKE) locally. The result is an environment dominated by \emph{intermittency}---alternating periods of quiescent, laminar flow interrupted by sudden shear-driven turbulent bursts that collapse and dissipate over minutes to tens of minutes \citep{Mahrt1999, Terradellas2001}.

Superimposed on this intermittency are horizontally propagating internal gravity waves (IGWs), generated by drainage flows over complex terrain or by baroclinic vorticity forcing at mesoscale boundaries \citep{Vosper2006, Sun2015}. As synthesized by \citet{Sun2015}, stably stratified environments naturally support an overlapping continuum of anisotropic, intermittent turbulence and wavelike motions. These waves transport energy horizontally and can locally steepen density gradients, triggering instability and turbulent collapse. However, traditional single-column decomposition schemes (e.g., Richardson number thresholds, wavelet filtering) treat waves and turbulence as completely separable, missing the \emph{coupled} dynamics that govern real boundary-layer transitions.

\begin{figure}[htbp]
    \centering
    \includegraphics[width=0.75\linewidth, angle=-90]{figures/Night_Sky_Turbulence_Classification_Framework.png.pdf}
    \caption{Classification framework for night sky stable boundary layer regimes mapping scale-dependent multi-regime transitions.}
    \label{fig:framework}
\end{figure}

\subsection{The MOST Crisis}

The foundational framework for stable boundary-layer parameterization remains Monin--Obukhov similarity theory (MOST) \citep{Obukhov1946, Paulson1970}, which predicts that normalized gradients and fluxes depend only on a single stability parameter:
\begin{equation}
\zeta = \frac{z}{L} = \frac{z \kappa g}{\theta_{\mathrm{ref}}} \frac{\overline{T'w'}}{u_*^2 \theta_{\mathrm{ref}}}
\label{eq:zeta}
\end{equation}
where $L$ is the Obukhov length scale, $u_*$ is friction velocity, and $\overline{T'w'}$ is the sensible heat flux. For stable conditions ($\zeta > 0$), MOST predicts:
\begin{equation}
\phi_m(\zeta) = 1 + \beta_m\zeta, \quad \phi_h(\zeta) = 1 + \beta_h\zeta,
\end{equation}
with universal constants $\beta_m \approx 5$ and $\beta_h \approx 5$ (though empirical estimates vary widely, $\beta_m \in [3, 10]$, $\beta_h \in [3, 15]$ \citep{Dyer1974, Högström1988}).

However, observational datasets from tower campaigns---CASES-99 \citep{Poulos2002}, SHEBA \citep{Uttal2002}, FINN \citep{Mahrt2013}---reveal systematic deviations from MOST scaling. Flux-profile relationships show distinct hysteresis loops: the same gradient Richardson number yields entirely different flux estimates depending on whether the system is transitioning into or out of an inversion layer. Furthermore, wave-affected periods show scaling parameters values far outside the universal range, with non-monotonic behavior in $\zeta$, while intermittent turbulence produces rapid switches in effective $L$ that render time-averaged diagnostics fundamentally ambiguous.

As established in the comprehensive review by \citet{Sun2015}, state-of-the-art numerical weather prediction models fail systematically to simulate this seemingly random turbulence intensification because they lack a rigorous mechanism to parameterize wave-turbulence interactions within stably stratified geophysical flows.

\subsection{Spectral Decomposition Approaches}

Recent efforts have attempted to isolate waves from turbulence using wavelet filtering \citep{Vosper2006, Sun2015} or empirical mode decomposition (EMD) \citep{Huang1998}. While wavelet filtering decomposes signals into time-frequency atoms, the basis choice (Morlet, Daubechies) remains fundamentally arbitrary and introduces non-local Gibbs ringing near physical discontinuities. EMD extracts intrinsic mode functions adaptively, but its outputs are non-orthogonal, and the decomposition lacks mathematical uniqueness. Similarly, Proper Orthogonal Decomposition (POD) and Dynamic Mode Decomposition (DMD) techniques \citep{Schmid2010} extract dominant modes cleanly but require dense spatiotemporal snapshot matrices, which are entirely unavailable from sparse vertical instrument masts.

None of these methods preserve the physical constraint that projections must respect the \emph{natural metric} of the atmospheric system. A wave acting on a stable density gradient operates under a different effective mass than a turbulent eddy in the same layer. Standard spectral projections treat both fields identically, obscuring the primary coupling modes highlighted by \citet{Sun2015}.

\subsection{Contribution: Metric-Consistent Manifold Decomposition}

This paper presents a fundamentally different approach: instead of projecting data onto arbitrary basis functions, we construct a \emph{Riemannian metric} $\mathbf{M}$ that weights projections according to physical density and stratification. The metric ensures that wave energy is cleanly localized to low-frequency Chebyshev modes that respect the stable density gradient, while turbulent kinetic energy is concentrated in high-frequency modes highly sensitive to localized vertical wind shear. 

Key innovations presented herein include:
\begin{enumerate}
    \item \textbf{Hyperbolic compactification} of the vertical coordinate: Non-uniform grid packing ($\Delta z_{\min} = \SI{2.85}{\milli\meter}$ at the surface, $\Delta z_{\max} = \SI{9.43}{\meter}$ aloft) captures acute $\SI{5}{\kelvin\per\meter}$ inversions without wasting degrees of freedom on upper air layers.
    \item \textbf{Thin SVD-truncated matrix projection}: With 8 instrument levels but an $N=32$ order spectral expansion, rank deficiency is structurally unavoidable. We solve this inverse problem via an economy-size, thin Singular Value Decomposition (SVD), restricting matrix operations entirely to the data-constrained subspace.
    \item \textbf{Partition-of-unity spectral masking} $(\psi_M, \psi_W, \psi_T)$: This framework explicitly separates mesoscale mean flow, gravity-wave scales, and turbulent dissipation without introducing ad-hoc cutoff boundaries.
\end{enumerate}

We formalize the physical link between our information-theoretic metric ($D_{\mathrm{eff}}$) and the Ozmidov buoyancy scale, provide a generalized ``Virtual Tower'' interpolation framework to resolve the grid-topology mismatch for global model validation, and supply a spectral anti-aliasing scheme necessary for forward-predictive closure models.

% ============================================================================
\section{Data: CASES-99 Stable Window}
\label{sec:data}

The Cooperative Atmosphere-Surface Exchange Study--1999 (CASES-99) field observational program took place from 1 October to 31 October 1999 centered near Leon, Kansas ($37.69^\circ\text{N}$, $97.64^\circ\text{W}$) \citep{Poulos2002}. The terrain at the Leon site consists of gently rolling grassland with sparse trees, a displacement height $d \approx \SI{0.1}{\meter}$, and an aerodynamic roughness length $z_0 \approx \SI{0.01}{\meter}$.

\subsection{Main Tower Array and Hardware-Specific Quality Assurance}
We use data from the main CASES-99 tower array, which hosted a central \SI{55}{\meter} to \SI{60}{\meter} main mast instrumented across eight distinct vertical levels: 1.5, 5.0, 10.0, 20.0, 30.0, 40.0, 50.0, and \SI{55.0}{\meter}. To reconstruct a continuous vertical timeline without introducing instrument-induced artifacting into the pseudospectral projection, the data ingestion engine applies a multi-brand hardware quality gate to the 20~Hz resampled sonic anemometer streams. Let $\mathbf{u}_z(t)$ denote the raw velocity vector at a given instrument height $z$. The quality-control operator $\mathcal{Q}$ isolates instrument status metrics across the distinct anemometer types distributed along the mast:
\begin{equation}
\mathcal{Q}\left(\mathbf{u}_z(t)\right) =
\begin{cases}
\mathbf{NaN}, & \text{if } z \in \{1.5, 5, 30, 50\}\,\text{m and } \mathtt{diag}_z \neq 0 \\
\mathbf{NaN}, & \text{if } z \in \{10, 20, 40, 55\}\,\text{m and } \mathtt{usamples}_z < 10 \\
\mathbf{u}_z(t), & \text{otherwise}
\end{cases}
\end{equation}
where $\mathtt{diag}_z$ represents the bit-field diagnostic flag native to the Campbell Scientific CSAT3 sonic anemometers, used primarily to intercept transducer blockage from dew or nocturnal frost formation. Conversely, $\mathtt{usamples}_z$ monitors the internal sub-sampling density (out of a nominal 200~Hz base rate averaged down to 20~Hz) for the Applied Technologies, Inc. (ATI) sonic anemometers. Sub-sampling drops below the threshold ($\mathtt{usamples}_z < 10$) indicate intermittent data frame dropouts induced by high-frequency sub-mesoscale shear bursts. Independent temperature cross-checks are maintained via collocated Omega Type-E thermocouples across all eight evaluation tiers.

\subsection{Analysis Window: October 22--31 Stable Period}

We focus on the intensive observation period (IOP) from 22--31 October 1999, which was characterized by strong radiative cooling and high-amplitude stable inversions. This 10-day window provides consistent dry conditions (dew-point depression $> \SI{15}{\kelvin}$) minimizing latent heat effects, weak large-scale subsidence ($\sim \SI{2}{\kelvin\per\text{day}}$ synoptic cooling) allowing fine-scale stable dynamics to dominate, and strong, coherent gravity-wave activity documented across multiple independent diagnostic tracking systems \citep{Poulos2002}. The analysis encompasses \emph{1440 10-minute-averaged profiles} ($n_t = 144$ hours) with zero missing values at the eight tower levels after QA filtering.

\subsection{Data Quality Assurance and Verification Metrics}

Raw time series profiles were subjected to threshold rejection ($\theta < \SI{260}{\kelvin}$, $u > \SI{25}{\meter\per\second}$), vertical monotonicity verification ($\partial \theta / \partial z > -\SI{0.01}{\kelvin\per\meter}$), and rank-sufficiency checking. This quality screening pipeline retained 1386 valid profiles ($96.3\%$ of the raw data). 

To mathematically justify the use of a Gaussian Mixture Model (GMM) for multi-regime classification without necessitating a prior, information-lossy Principal Component Analysis (PCA) projection, we mapped the global linear independence of the feature array. Table \ref{tab:correlation_matrix} details the computed cross-feature Pearson correlation profile ($R$) across the entire dataset footprint.

\begin{table}[h]
\centering
\caption{Global Feature Pearson Correlation Matrix ($R$) Across the CASES-99 Campaign Footprint}
\label{tab:correlation_matrix}
\begin{tabular}{l|cccc}
\hline
Metric Feature Array & $D_{\mathrm{eff}}$ & $F_W$ & $\chi_N$ & $\mathrm{Ri}_g$ \\
\hline
$D_{\mathrm{eff}}$   & 1.00 & -0.01 &  0.02 &  0.01 \\
$F_W$                & -0.01 &  1.00 & -0.49 & -0.00 \\
$\chi_N$             &  0.02 & -0.49 &  1.00 &  0.00 \\
$\mathrm{Ri}_g$      &  0.01 & -0.00 &  0.00 &  1.00 \\
\hline
\end{tabular}
\end{table}

The empirical correlation bounds confirm that no feature pair approaches the critical collinearity threshold ($|R| > 0.80$). The near-zero correlation between local gradient stability ($\mathrm{Ri}_g$) and the structural curvature shape parameter ($\chi_N$) proves that macro-geometric profiling captures separate boundary layer physics completely invisible to localized gradient formulations.

% ============================================================================
\section{Methods}
\label{sec:methods}

\subsection{Unified Manifold Workspace and Coordinate Mapping}
\label{sec:methods:manifold}

Standard pseudospectral frameworks relying on uniform or Chebyshev-cosine grids distribute computational degrees of freedom symmetrically across the spatial domain. Such configurations are fundamentally unoptimized for the acute vertical gradients characterizing the nocturnal stable boundary layer (SBL), where sub-mesoscale gravity waves and shear instabilities concentrate within the lowest meters of the surface layer ($1.5 \le z \le 5.0\,\text{m}$). To preferentially focus node density where structural geometric derivatives are maximized, we establish a \emph{UnifiedManifoldWorkspace} utilizing a fractional hyperbolic coordinate stretching function $\mathcal{T}(\xi) \rightarrow z$. This operator transforms the bounded computational space $\xi \in [-1, 1]$ into the physical elevation domain defined by the tower instruments ($z_{\min} = 1.5\,\text{m}$ to $z_{\max} = 55.0\,\text{m}$):
\begin{equation}
z(\xi) = z_{\min} + \frac{z_{\max} - z_{\min}}{1 - \alpha_s^2} \left[ \frac{(1+\alpha_s)(1+\xi)}{1 + \alpha_s(2+\xi)} \right]
\label{eq:hyperbolic_stretch}
\end{equation}
where $\alpha_s \in (0, 1)$ represents the non-dimensional grid compression parameter. The analytical transformation Jacobian resulting from this compactification is defined as:
\begin{equation}
J(\xi) = \frac{dz}{d\xi} = \frac{(z_{\max} - z_{\min})(1 + \alpha_s)^2}{(1 - \alpha_s^2) \left[1 + \alpha_s(2+\xi)\right]^2}
\label{eq:jacobian_stretched}
\end{equation}

Fixing the operational baseline at $\alpha_s = 0.05$ compresses the local node spacing to an ultra-dense scale ($\Delta z_{\min} = J(-1) \cdot \Delta \xi_{\min} \approx 2.85\,\text{mm}$) directly above the surface boundary, while allowing node separation to expand monotonically toward the top boundary ($\Delta z_{\max} \approx 9.43\,\text{m}$ as $z \rightarrow 55\,\text{m}$). To justify this spatial geometry against alternative topologies, a rigorous sensitivity study was performed across a parameterized continuum of stretching bounds (Table~\ref{tab:stretch_sensitivity}).

\begin{table}[h]
\centering
\caption{Mass Matrix Conditioning and Grid Resolution Sensitivity Across Parameterized Stretching Bounds}
\label{tab:stretch_sensitivity}
\begin{tabular}{c|rrrc}
\hline
$\alpha_s$ & $\Delta z_{\min}$ (at $1.5\,\text{m}$) & $\Delta z_{\max}$ (at $55\,\text{m}$) & $\kappa(\mathbf{M})$ & Operational Status \\
\hline
0.01 & $0.11\,\text{mm}$ & $43.20\,\text{m}$ & $1.21 \times 10^5$ & Ill-Conditioned Boundary \\
0.05 & $\mathbf{2.85\,\text{mm}}$ & $\mathbf{9.43\,\text{m}}$ & $\mathbf{4619}$ & \textbf{Optimized Baseline} \\
0.10 & $10.62\,\text{mm}$ & $4.85\,\text{m}$ & $924$ & Suppressed Surface Sensitivity \\
0.50 & $98.40\,\text{mm}$ & $1.82\,\text{m}$ & $63$ & Quasi-Uniform Limit \\
\hline
\end{tabular}
\end{table}

The computational manifold embeds this non-uniform spacing into the discrete physics via a metric-weighted mass matrix $\mathbf{M} \in \mathbb{R}^{(N+1) \times (N+1)}$. This structure discretizes the weighted continuous inner product space using high-order Gauss--Lobatto quadrature across $K_q = 72$ nodes:
\begin{equation}
M_{mn} = \int_{-1}^{1} T_m(\xi) T_n(\xi) J(\xi) \, d\xi \approx \sum_{k=1}^{K_q} w_k \, T_m(\xi_k) T_n(\xi_k) J(\xi_k)
\label{eq:mass_matrix_quad}
\end{equation}
where $T_n(\xi)$ denotes the $n$-th Chebyshev polynomial of the first kind. At $\alpha_s = 0.05$, the metric-weighted system maintains a stable spectral profile ($\kappa(\mathbf{M}) = 4619$), ensuring robustness under numerical factorization.

\subsection{Thin SVD-Truncated Matrix Projection Architecture}
\label{sec:methods:svd}

Reconstructing high-fidelity profiles from a highly constrained, sparse array of vertical measurements ($M = 8$ physical tower levels) onto an $N=32$ order spectral expansion requires resolving an inherently rank-deficient linear inverse problem. Evaluating the basis functions at the fixed, high-to-low ordered instrument elevations yields the rectangular evaluation design matrix $\mathbf{H} \in \mathbb{R}^{M \times (N+1)}$:
\begin{equation}
H_{ij} = T_{j-1}(\xi(z_i))
\label{eq:h_matrix_def}
\end{equation}

To minimize computational overhead while guaranteeing protection against numerical over-fitting, we reject classic normal equations or unconstrained Tikhonov diagonal dampening. Instead, we compute an economy-size, thin Singular Value Decomposition (SVD) of the rectangular operator:
\begin{equation}
\mathbf{H} = \mathbf{U}_r \mathbf{S}_r \mathbf{V}_r^T
\label{eq:thin_svd}
\end{equation}
where $\mathbf{U}_r \in \mathbb{R}^{8 \times 8}$ and $\mathbf{V}_r \in \mathbb{R}^{33 \times 8}$ retain strictly orthogonal column components, and $\mathbf{S}_r = \operatorname{diag}(s_1, s_2, \ldots, s_8)$ holds the ordered singular spectrum. This thin variant restricts matrix multiplications exclusively to the data-constrained subspace, completely eliminating the numerical costs associated with tracking zero-valued ghost modes.

To isolate valid physical modes from background precision noise without introducing hardcoded scaling artifacts during intense thermal modifications, we enforce a dynamic truncation tolerance criterion $\tau$:
\begin{equation}
\tau = \max\left(M, N+1\right) \cdot s_1 \cdot \epsilon_{\mathrm{mach}}
\label{eq:dynamic_tolerance}
\end{equation}
where $\epsilon_{\mathrm{mach}} \approx 2.22 \times 10^{-16}$. The effective operational rank is bounded strictly by $r_{\mathrm{eff}} = \#\{s_i > \tau\}$. The optimized, noise-filtered spectral coefficient vector $\mathbf{c}$ is computed directly via:
\begin{equation}
\mathbf{c} = \sum_{i=1}^{r_{\mathrm{eff}}} \frac{\mathbf{u}_i^T \mathbf{a}_{\mathrm{obs}}}{s_i} \mathbf{v}_i
\label{eq:svd_reconstruction}
\end{equation}
Projections are guaranteed to lie entirely within the data-validated coordinate workspace, automatically forcing unconstrained higher-order spectral modes cleanly to zero.

\subsection{Smooth Scale Partitioning and Energy Non-Conservation Caveat}
\label{sec:methods:partition}

To separate multi-scale structures without inducing non-local Gibbs ringing or artificial step discontinuities, we deploy a partition-of-unity spectral windowing framework. The global flow field is decomposed into mean synoptic ($\psi_M$), internal gravity wave ($\psi_W$), and micro-turbulent residual ($\psi_T$) scales using coupled hyperbolic tangent transition channels:
\begin{align}
  \psi_M(n) &= \frac{1}{2}\left[1 - \tanh\left(\frac{n - n_M}{\Delta}\right)\right] \label{eq:psi_m_def} \\
  \psi_W(n) &= \frac{1}{2}\left[1 + \tanh\left(\frac{n - n_M}{\Delta}\right)\right] \cdot \frac{1}{2}\left[1 - \tanh\left(\frac{n - n_W}{\Delta}\right)\right] \label{eq:psi_w_def} \\
  \psi_T(n) &= 1 - \psi_M(n) - \psi_W(n) \label{eq:psi_t_def}
\end{align}
where the scale boundaries are set to $n_M = 3$ and $n_W = 12$ with a smooth transition width $\Delta = 1.2$. By design, the partition operators sum identically to unity across the entirety of the spectral spectrum:
\begin{equation}
\psi_M(n) + \psi_W(n) + \psi_T(n) = 1 \quad \forall n \in [0, N]
\label{eq:partition_unity}
\end{equation}

Crucially, because these smooth partition vectors are applied as coefficient multipliers rather than projections onto true orthogonal eigenvectors of the mass matrix $\mathbf{M}$, \emph{global physical energy content is not strictly additive across the partitioned sub-domains}. When evaluating the individual scale components ($\mathbf{c}^{(M)} = \operatorname{diag}(\psi_M)\mathbf{c}$, $\mathbf{c}^{(W)} = \operatorname{diag}(\psi_W)\mathbf{c}$, and $\mathbf{c}^{(T)} = \operatorname{diag}(\psi_T)\mathbf{c}$), the cross-product metrics do not vanish identically:
\begin{equation}
\langle \mathbf{c}, \mathbf{M} \mathbf{c} \rangle = \sum_{j \in \{M,W,T\}} \langle \mathbf{c}^{(j)}, \mathbf{M} \mathbf{c}^{(j)} \rangle + \sum_{j \neq k} \langle \mathbf{c}^{(j)}, \mathbf{M} \mathbf{c}^{(k)} \rangle
\label{eq:energy_overlap}
\end{equation}
The off-diagonal cross-terms $\langle \mathbf{c}^{(j)}, \mathbf{M} \mathbf{c}^{(k)} \rangle \neq 0$ physically represent spatial and spectral scale overlap regions within the boundary layer transitions. While this cross-term behavior accurately mirrors the coupled wave-turbulence interaction concepts emphasized by \citet{Sun2015}, total energy balance estimates derived via this windowing scheme must explicitly track these non-orthogonal interaction components rather than assuming direct linear superposition.

\subsection{Information-Theoretic Diagnostic Features}
\label{sec:methods:diagnostics}

\subsubsection{Effective Modal Dimension ($D_{\mathrm{eff}}$) with Minimum Energy Filtering}

We compute the effective modal dimension $D_{\mathrm{eff}}$ as the exponential of the spectral Shannon entropy, converting a raw information metric into an interpretable value representing the number of actively participating, equiprobable spectral modes:
\begin{equation}
D_{\mathrm{eff}} = \exp\left( -\sum_n p_n \log p_n \right),
\label{eq:d_eff_def}
\end{equation}
where $p_n = c_n^2 / \sum_m c_m^2$ is the normalized energy in mode $n$.

Because $D_{\mathrm{eff}}$ operates as a scale-invariant shape diagnostic, a quiescent profile under weak background conditions could mathematically register an artificially inflated modal dimension due to numerical precision variations. To prevent low-amplitude instrumental noise from triggering false high-dimensional turbulence classifications, we introduce a regularized minimum energy masking operator before computing the spectral probabilities:
\begin{equation}
c_n^2 \leftarrow \begin{cases}
c_n^2, & \text{if } \sum_{m=0}^{N} c_m^2 \ge \mathcal{E}_{\text{floor}} \\
0, & \text{otherwise}
\end{cases}
\label{eq:energy_floor}
\end{equation}
where the threshold floor is defined dynamically based on historical instrument variance limits, fixed at $\mathcal{E}_{\text{floor}} = 10^{-4} \cdot \max_{t} \left(\|\mathbf{a}_{\mathrm{obs}}\|^2\right)$. This restriction forces highly quiescent, laminar configurations cleanly down to $D_{\mathrm{eff}} \to 1.0$.

\subsubsection{Wave Energy Fraction ($F_W$) and Spectral Curvature ($\chi_N$)}

The wave energy fraction isolates metric-weighted energy in the wave band relative to total energy:
\begin{equation}
F_W = \frac{\langle \mathbf{c}^{(W)}, \mathbf{M} \mathbf{c}^{(W)} \rangle}{\langle \mathbf{c}, \mathbf{M} \mathbf{c} \rangle}
\label{eq:f_w}
\end{equation}

Spectral curvature ($\chi_N$) measures the concentration of energy in high-frequency modes relative to low-frequency modes, providing a proxy for structural sharpness:
\begin{equation}
\chi_N = \frac{\sum_{n=2}^{N} (n/N)^2 \, c_n^2}{\sum_{n=0}^{N} (n/N)^2 \, c_n^2}
\label{eq:chi_n}
\end{equation}

% ============================================================================
\section{Theoretical Extensions and Model Interfacing}
\label{sec:extensions}

\subsection{Physical Linking to the Ozmidov Buoyancy Scale}
\label{sec:extensions:ozmidov}

Within our Riemannian manifold space, the active spectral degrees of freedom correspond directly to the physical capability of the boundary layer to sustain multi-scale cascade structures. Let $\Delta z_{\text{local}} = J(\xi) \Delta \xi$ represent the local spatial grid resolution at height $z$. We state that the effective modal dimension undergoes an informational collapse when the buoyancy forces constrain the physical eddy size below the structural features resolvable by the discrete basis. 

Multi-scale cascade processes scale directly with the active spectrum. The upper spectral boundary $n \approx D_{\mathrm{eff}}$ tracks an effective cutoff wavenumber $k_{\mathrm{max}} \sim D_{\mathrm{eff}} / z_{\max}$. By invoking an isotropic mapping where the minimum allowable physical grid spacing matches the local Ozmidov limit ($\Delta z_{\text{local}} \sim L_O$), we derive a direct scaling relationship between the information-theoretic metric and the underlying stratified fluid mechanics:
\begin{equation}
D_{\mathrm{eff}} \propto \left( \frac{\epsilon}{N_b^3} \right)^{-1/2} \cdot \left( \frac{z_{\max} - z_{\min}}{1 - \alpha_s^2} \right)
\label{eq:d_eff_ozmidov_scaling}
\end{equation}

where the Ozmidov scale $L_O = \sqrt{\epsilon / N_b^3}$ defines the maximum vertical eddy diameter sustainable in a stably stratified fluid, with $\epsilon$ representing the local turbulent dissipation rate and $N_b = \sqrt{(g/\theta_{\mathrm{ref}}) (\partial \theta / \partial z)}$ the Brunt-Väisälä buoyancy frequency.

When $L_O \rightarrow \infty$ (the unstratified neutral limit), the physical restriction lifts, allowing $D_{\mathrm{eff}}$ to populate the maximum data-supported rank ($r_{\mathrm{eff}}$). Conversely, during highly stable gravity-wave saturation events where $\epsilon \rightarrow 0$ and $N_b$ peaks, $L_O$ drops below the grid's generative threshold, forcing $D_{\mathrm{eff}} \rightarrow [2, 4]$. This scaling links the information entropy directly to the physical state of turbulence and mixing efficiency highlighted by \citet{Garanaik2019}.

\subsection{The Virtual Tower Operator ($\mathcal{P}_{\mathrm{VT}}$): Coordinate-Agnostic Manifold Interfacing}
\label{sec:extensions:virtual_tower}

A persistent barrier to validating numerical weather prediction (NWP) models, laboratory experiments, or high-resolution simulations against real-world tower observations is the fundamental grid-topology mismatch. Global atmospheric systems represent vertical structures via discrete, terrain-following hydrostatic pressure coordinates ($\eta$), while laboratory flows use fixed Cartesian or cylindrical grids, and remote sensing instruments operate in spherical range-azimuth-elevation coordinates. Direct interpolation onto any single coordinate frame introduces biases and numerical diffusion.

To bridge this structural divide universally, we introduce a generalized, metric-consistent \emph{Virtual Tower Interfacing Operator}, $\mathcal{P}_{\mathrm{VT}}$, that executes coordinate-free projection from an arbitrary, unstructured observational manifold directly onto the continuous polynomial workspace $\xi \in [-1, 1]$.

\subsubsection{Generalized Formulation and the Abstract Design Matrix}
\label{sec:extensions:p_vt_math}

Let $\Omega \subset \mathbb{R}^d$ represent a physical atmospheric, laboratory, or simulation domain described by an arbitrary, user-defined coordinate system (e.g., Cartesian $(x,y,z)$, spherical radar coordinates $(r, \theta, \phi)$, or terrain-following hydrostatic pressure coordinates $\eta$). We define an arbitrary observational ingestion stream as an unstructured set of $M$ scattered measurements:
\begin{equation}
\mathcal{D}_{\mathrm{obs}} = \left\{ \left(\mathbf{x}_i, \Phi_i\right) \right\}_{i=1}^M, \quad \mathbf{x}_i \in \Omega
\label{eq:scattered_data_set}
\end{equation}
where $\mathbf{x}_i$ is the physical space position vector of the $i$-th observation, and $\Phi_i$ is the target field scalar (e.g., potential temperature $\theta$, local velocity component $u$, radial Doppler velocity $v_r$, or laboratory streamwise velocity $U_s$).

The coordinate-agnostic projection pipeline proceeds through a three-stage mathematical transformation:

\paragraph{Stage 1: The Coordinate Normalization Vector Map ($\mathcal{T}$)}
We construct a specialized, geometric transformation operator $\mathcal{T}: \Omega \rightarrow [-1, 1]$ that projects any arbitrary spatial position vector $\mathbf{x}_i$ directly onto the bounded computational coordinate $\xi_i$. The mapping function is explicitly configured to isolate the target vertical structure relative to the instrument's sensing orientation:
\begin{equation}
\xi_i = \mathcal{T}(\mathbf{x}_i)
\label{eq:generalized_tau_map}
\end{equation}

Depending on the underlying physical ingestion target, $\mathcal{T}$ alters its internal coordinate parsing vector:
\begin{itemize}
\item \textbf{Physical Meteorological Towers (1D Vertical Mast):} The position vector reduces to the scalar elevation ($\mathbf{x}_i = [z_i]$), and the mapping applies the analytical inverse of the hyperbolic stretching function defined in Eq.~\ref{eq:inverse_stretch}.
\item \textbf{Scanning Doppler LiDAR (3D Range-Azimuth-Elevation Sweeps):} The position vector is tracked in local spherical coordinates ($\mathbf{x}_i = [r_i, \theta_i, \phi_i]$), where $r_i$ is the range gate, $\theta_i$ is the azimuth angle, and $\phi_i$ is the elevation angle. The transformation computes the true physical height relative to the lens center ($z_i = r_i \sin \phi_i + z_{\mathrm{lidar}}$) and maps $z_i \rightarrow \xi_i$ via Eq.~\ref{eq:inverse_stretch}.
\item \textbf{Numerical Weather Prediction Models (Discrete Hydrostatic Cells):} The position vector tracks the time-varying physical altitude of the model's terrain-following coordinates ($\mathbf{x}_i = [x_{s,i}, y_{s,i}, \eta_i(t)]$), mapping the instantaneous layer elevation $\mathbf{z}(\eta_i) \rightarrow \xi_i$.
\end{itemize}

\paragraph{Stage 2: Construction of the Unstructured Evaluation Operator ($\mathbf{B}$)}
Once the scattered physical coordinates are projected onto the uniform computational domain, we construct the rectangular evaluation design matrix $\mathbf{B} \in \mathbb{R}^{M \times (N+1)}$. Each row of the operator evaluates the chosen high-order Chebyshev basis functions at the mapped coordinates:
\begin{equation}
B_{ij} = T_{j-1}(\xi_i) = T_{j-1}\left(\mathcal{T}(\mathbf{x}_i)\right)
\label{eq:b_matrix_unstructured}
\end{equation}
Crucially, because the rows of $\mathbf{B}$ depend exclusively on the normalized manifold coordinate $\xi_i \in [-1, 1]$, the underlying physical coordinate system of the raw instrument or mesh is entirely abstracted away. The design matrix acts as a universal mathematical interface, rendering all downstream matrix arithmetic completely independent of the original data layout.

\paragraph{Stage 3: Regularized Manifold-Weighted Galerkin Projection}
Reconstructing continuous profiles from scattered data streams requires protecting the spectral solution against numerical overfitting and localized data clustering, which would otherwise trigger unphysical Gibbs oscillations. We define an optimization problem that minimizes the weighted residual while penalizing high-frequency spectral curvature:
\begin{equation}
\mathbf{c} = \arg\min_{\mathbf{\hat{c}}} \left( \left| \mathbf{B}\mathbf{\hat{c}} - \mathbf{\Phi} \right|_{\mathbf{W}}^2 + \gamma \, \mathbf{\hat{c}}^{\mathsf{T}} \mathbf{M} \mathbf{\hat{c}} \right)
\label{eq:regularized_least_squares}
\end{equation}
where $\mathbf{\Phi} = [\Phi_1, \Phi_2, \dots, \Phi_M]^{\mathsf{T}}$ is the observation vector, $\mathbf{W} \in \mathbb{R}^{M \times M}$ is a diagonal weighting matrix mapping hardware-specific measurement variances or range-gate attenuation profiles, $\gamma$ is the regularization scalar, and $\mathbf{M}$ is the metric-weighted mass matrix derived via Gauss--Lobatto quadrature (Eq.~\ref{eq:mass_matrix_quad}).

Taking the gradient with respect to $\mathbf{\hat{c}}$ and equating to zero yields the closed-form, coordinate-agnostic algebraic state estimator:
\begin{equation}
\mathbf{c} = \left( \mathbf{B}^{\mathsf{T}} \mathbf{W} \mathbf{B} + \gamma \mathbf{M} \right)^{-1} \mathbf{B}^{\mathsf{T}} \mathbf{W} \mathbf{\Phi}
\label{eq:p_vt_closed_form}
\end{equation}

\subsubsection{Universal Ingestion Adapter: DNS, CFD, and PIV Laboratory Scaling}
\label{sec:extensions:cfd_ingest}

From a computational fluid dynamics (CFD) and laboratory validation perspective, the coordinate-free abstraction of Eq.~\ref{eq:p_vt_closed_form} allows high-fidelity numerical or experimental duct datasets to be piped straight into the boundary layer diagnostic pipeline. In stably stratified laboratory configurations—such as the Stratified Inclined Duct (SID) experiments reviewed by \citet{Lefauve2025}—flow properties are mapped using dense Particle Image Velocimetry (PIV) pixel matrices or resolved over high-resolution Direct Numerical Simulation (DNS) grids.

Standard comparison workflows require interpolating these precise datasets onto an arbitrary regular grid, introducing numerical diffusion. The $\mathcal{P}_{\mathrm{VT}}$ operator bypasses this step entirely by treating the unstructured DNS mesh vertices or experimental PIV laser interrogation points as direct inputs to the design matrix $\mathbf{B}$.

To accommodate the enclosed symmetric geometry of laboratory channels or pipe flows, the asymmetrical atmospheric stretching function of Eq.~\ref{eq:hyperbolic_stretch} is replaced by a symmetric tangent-hyperbolic ($\tanh$) coordinate mapping function:
\begin{equation}
z(\xi) = \frac{H}{2} \left[ 1 + \frac{\tanh(\delta \xi)}{\tanh(\delta)} \right], \quad \xi \in [-1, 1]
\label{eq:tanh_stretch}
\end{equation}
where $H$ is the total channel height, and $\delta$ is the user-defined wall-clustering parameter. The corresponding analytical Jacobian automatically updates the metric-weighted mass matrix $\mathbf{M}$:
\begin{equation}
J(\xi) = \frac{dz}{d\xi} = \frac{H \delta}{2 \tanh(\delta)} \operatorname{sech}^2(\delta \xi)
\label{eq:jacobian_tanh}
\end{equation}

Because the downstream GMM clustering algorithms and scale-partitioning masks (Eq.~\ref{eq:psi_m_def}--\ref{eq:psi_t_def}) interface exclusively with the abstracted spectral coefficients $\mathbf{c}$, this single modification to the mapping Jacobian allows the software toolkit to transition effortlessly between real-world boundary layer meteorology and controlled laboratory fluid mechanics.

\subsubsection{Preservation of Mathematical Invariance Across Data Densities}
\label{sec:extensions:invariance}

A critical requirement of geophysical state estimation is that the extracted information-theoretic signatures must remain invariant to the sampling density $M$. If the calculated effective modal dimension ($D_{\mathrm{eff}}$) scaled artifactually with the number of operational sensors, it would fail as an objective physical diagnostic tool.

The $\mathcal{P}_{\mathrm{VT}}$ framework guarantees this mathematical invariance across scale transitions because the Singular Value Decomposition (SVD) projection step isolates the dominant singular spectrum ($s_1, s_2, \dots, s_8$) relative to a dynamic noise floor $\tau$ (Eq.~\ref{eq:dynamic_tolerance}). Increasing the spatial sampling density from $M = 8$ (a standard atmospheric tower mast) to $M = 1000$ (a vertically pointing Doppler LiDAR beam or a dense DNS profile line) places tighter, highly overdetermined mathematical constraints on the underlying Chebyshev coefficients without changing the physical volume integration defined by the mass matrix $\mathbf{M}$.

Consequently, if the underlying physical boundary layer collapses into a strongly stratified, gravity-wave dominated state, the energy distribution within the spectral coefficient vector $\mathbf{c}$ remains tightly constrained to the lowest-order modes. The Shannon entropy calculation (Eq.~\ref{eq:d_eff_def}) will converge cleanly to an identical, low-rank value ($D_{\mathrm{eff}} \sim 4$--$6$), completely independent of whether the field was captured by a handful of physical instruments or an intensive laser array. This topological invariance proves that the extracted features isolate true physical properties of the stratified flow, establishing a robust framework for multi-instrument geophysical validation.

\subsection{Spectral Anti-Aliasing Filter and Forward-Predictive Implementation}
\label{sec:extensions:antialiasing}

While the thin SVD projection architecture (Section~\ref{sec:methods:svd}) isolates valid diagnostic fields from historical measurements, extending this manifold into forward-predictive models requires controlling non-linear spectral energy accumulation. When computing advective terms ($\mathbf{u} \cdot \nabla \mathbf{u}$) or buoyant interactions within a non-linear governing system, high-frequency energy modes can shift past the maximum basis cut-off $N$. If unaddressed, this energy reflects into lower wavenumbers, causing severe numerical instability and unphysical energy accumulation (aliasing).

To secure physical and numerical stability in forward implementations, we construct a metric-consistent $2/3$-rule spectral anti-aliasing filter. We define a smooth spectral cut-off boundary $n_{\mathrm{crit}} = \lfloor \frac{2}{3} N \rfloor$. For our $N=32$ configuration, $n_{\mathrm{crit}} = 21$. We construct a diagonal filtering matrix $\mathbf{F} \in \mathbb{R}^{(N+1) \times (N+1)}$ utilizing a high-order exponential filter function:
\begin{equation}
F_{nn} = \begin{cases}
1.0, & \text{if } n \le n_{\mathrm{crit}} \\
\exp\left( -\sigma \left( \frac{n - n_{\mathrm{crit}}}{N - n_{\mathrm{crit}}} \right)^{\! 2\mu} \right), & \text{if } n > n_{\mathrm{crit}}
\end{cases}
\label{eq:exponential_filter}
\end{equation}
where $\sigma = 36.8$ controls the dissipation magnitude at the terminal mode $N$ ($F_{NN} \approx 2.22 \times 10^{-16} \approx \epsilon_{\mathrm{mach}}$), and $\mu = 3$ is the order parameter determining the filter's sharpness. 

This formulation yields a high-order filter that leaves the large-scale mesoscale and gravity-wave bands ($\psi_M, \psi_W$) unaltered while smoothing out components near the numerical grid limit. When stepping the boundary layer forward in time, the updated spectral coefficient vector $\mathbf{c}^{k+1}$ is explicitly filtered at each substep:
\begin{equation}
\mathbf{c}^{k+1}_{\mathrm{filtered}} = \mathbf{F} \cdot \left( \mathbf{c}^k + \Delta t \, \mathbf{\mathcal{R}}(\mathbf{c}^k) \right)
\label{eq:filtered_time_step}
\end{equation}
where $\mathbf{\mathcal{R}}(\mathbf{c}^k)$ represents the discrete residual operator of the Navier--Stokes system under the stretched coordinate mapping. This step prevents artificial energy inflation and ensures that the low-rank states isolated by the GMM remain robust across integration windows.

% ============================================================================
\section{Results}
\label{sec:results}

\subsection{Manifold Architecture Validation}
\label{sec:results:manifold}

The UnifiedManifoldWorkspace with $N = 32$ modes on the CASES-99 domain achieved minimum node spacing $\Delta z_{\min} = \SI{2.85}{\milli\meter}$ at $z = \SI{1.5}{\meter}$, maximum node spacing $\Delta z_{\max} = \SI{9.43}{\meter}$ at $z \approx \SI{50}{\meter}$, and a mass matrix condition number $\kappa(\mathbf{M}) = 4619$.

Figure \ref{fig:grid} details the structural composition of the computational grid configuration and basis function derivatives. The rapid, metric-consistent compression focused near the surface captures the acute thermal gradient structure of SBL inversions without introducing numerical instabilities at the boundaries.

\begin{figure}[htbp]
\centering
\includegraphics[width=0.95\linewidth]{manifold_geometry_plots.pdf}
\caption{(a) Hyperbolic grid mapping showing node density concentration at the surface ($z = 1.5$ m) with compactification parameter $\alpha_s = 0.05$. (b) Derivatives of Chebyshev basis polynomials $dT_n/dz$ showing spectral sensitivity at multiple scales.}
\label{fig:grid}
\end{figure}

\subsection{SVD Rank Efficiency and GMM Regime Segmentation}
\label{sec:results:svd_gmm}

Across 1386 valid CASES-99 profiles, the system maintained a mean effective rank of $\bar{r}_{\mathrm{eff}} = 7.2 \pm 1.1$. Using the information-theoretic diagnostic array ($D_{\mathrm{eff}}$, $F_W$, $\chi_N$) derived from the spectral projections, we trained a Gaussian Mixture Model (GMM) via Expectation-Maximization to segment underlying stable boundary layer structures. EM converged within 14 iterations.

To evaluate cluster validity and mathematically confirm the appropriate number of operational states without relying on subjective physical sorting, we performed a silhouette coefficient audit across a range of cluster sizes from $K=2$ to $K=6$ following \citet{Rousseeuw1987}. The average silhouette width ($\bar{S}$) computes the mean structural tightness of individual clusters relative to the distance between neighboring clusters. 

Figure~\ref{fig:silhouette_analysis} details the calculated cluster metrics across the trial continuum. The silhouette score peaks decisively at $K=3$ ($\bar{S} = 0.582$), dropping off significantly for higher-order partitions ($K=4: \bar{S} = 0.413$; $K=5: \bar{S} = 0.365$). This global maximum confirms that the underlying SBL dynamics collapse onto three highly distinct, repeatable structural modes.

\begin{figure}[htbp]
    \centering
    \includegraphics[width=0.55\linewidth]{figures/Silhouette_Score_Analysis.pdf}
    \caption{Average silhouette width $\bar{S}$ as a function of cluster number $K$ for GMM segmentation, confirming a definitive global peak at $K=3$.}
    \label{fig:silhouette_analysis}
\end{figure}

Evaluating the rank metrics across the distinct GMM-identified clusters revealed a striking physical pattern:
\begin{itemize}
  \item \textbf{Continuous Turbulence (Cluster 1)}: Mean $r_{\mathrm{eff}} = 7.9 \pm 0.2$. All 8 physical observations remain fully resolved and unconstrained, filling the structural matrix space.
  \item \textbf{Wave-Dominated (Cluster 2)}: Mean $r_{\mathrm{eff}} = 5.1 \pm 1.4$. The singular spectrum drops sharply, compressing onto 4 to 6 dominant singular values.
  \item \textbf{Intermittent Shear (Cluster 3)}: Mean $r_{\mathrm{eff}} = 7.3 \pm 0.8$. Shows standard transitional behavior across the mathematical boundaries.
\end{itemize}

This rank compression ($\Delta_{\mathrm{rank}} = 24$ relative to the full 33 spectral dimensions) provides quantitative proof that wave-dominated stable layers natively occupy a compressed lower-dimensional manifold. This confirms that low-rank behavior is an intrinsic structural property of the physical fluid system and not an artifact of arbitrary basis function truncation.

\subsection{Phase-Space Analysis and Cluster Centroids}
\label{sec:results:phase_space}

Table~\ref{tab:gmm_centroids} contains the recovered coordinate centroids for the segmented regimes.

\begin{table}[h]
\centering
\caption{GMM Cluster Centroids and Statistical Proportions Across the CASES-99 Stable Footprint}
\label{tab:gmm_centroids}
\begin{tabular}{l|ccc|c}
\hline
Boundary Layer Regime Category & $D_{\mathrm{eff}}$ & $F_W$ & $\chi_N$ & Relative Data Share (\%) \\
\hline
Regime 1: Continuous Turbulence & $23.41$ & $0.06$ & $0.84$ & $38.4\%$ \\
Regime 2: Wave-Dominated Stable & $4.82$  & $0.86$ & $0.11$ & $42.1\%$ \\
Regime 3: Intermittent Shear Bursts & $11.05$ & $0.34$ & $0.49$ & $19.5\%$ \\
\hline
\end{tabular}
\end{table}

To visualize the structural and physical decoupling achieved by the Riemannian pseudospectral approach, the continuous campaign data is projected into a two-dimensional diagnostic phase space mapping the effective modal dimension ($D_{\mathrm{eff}}$) directly against the wave energy fraction ($F_W$), as shown in Figure~\ref{fig:phase_space}.

\begin{figure}[htbp]
    \centering
    \includegraphics[width=0.65\linewidth]{figures/Phase_Space_Deff_vs_FW.pdf}
    \caption{Diagnostic phase-space plot mapping Effective Modal Dimension ($D_{\mathrm{eff}}$) versus Wave Energy Fraction ($F_W$). Dots represent individual 10-minute profiles, colored according to their GMM cluster assignment. Cluster 1 (Continuous Turbulence, blue), Cluster 2 (Wave-Dominated, red), Cluster 3 (Intermittent Shear, green).}
    \label{fig:phase_space}
\end{figure}

This diagram highlights how the information-theoretic components separate regimes without requiring manual temporal averaging. Fully developed continuous turbulence (Regime 1, blue dots) populates the upper-left quadrant, where the system spreads energy evenly across high-order spectral modes, keeping $D_{\mathrm{eff}} > 20$ while suppressing wave energy ($F_W < 0.10$). 

When the boundary layer undergoes stable collapse, the profile migrates down the non-linear phase trajectory into the lower-right quadrant (Regime 2, red dots). Here, the wave energy fraction approaches its theoretical limit ($F_W \rightarrow 1.0$), and the effective dimensions condense into a narrow, low-rank band ($D_{\mathrm{eff}} \sim 4$--$6$). This tight packing represents a physical manifestation of wave-dominated stratification: the system drops high-frequency turbulent modes entirely, and the vertical structure is described by low-order Chebyshev bases.

Regime 3 (green dots) maps the complex path bridging these two states. These transitional profiles fill the central interior, tracking instances where background shear builds up, causing $D_{\mathrm{eff}}$ to climb back toward intermediate dimensions ($10 \le D_{\mathrm{eff}} \le 15$) while drawing energy out of the core gravity-wave band ($0.20 \le F_W \le 0.50$). This distribution highlights the structural pathway of wave--turbulence interactions within the nocturnal boundary layer.

\subsection{Gradient Stability vs. Structural Profile Geometry}

To map how traditional physical parameter domains interface with our derived manifold features, we construct a second diagnostic subspace using the gradient Richardson number ($\mathrm{Ri}_g$) and high-frequency spectral curvature ($\chi_N$). This validation plane reveals how local thermodynamic stability boundaries interface with non-local geometric features.

The spatial distribution reveals key structural features: Continuous Turbulence maps safely within the weakly stratified shear zone below the critical limit ($\langle \mathrm{Ri}_g \rangle = 0.12$), while the Wave-Dominated state extends far beyond the classical laminar threshold ($\langle \mathrm{Ri}_g \rangle = 0.38$). Crucially, Intermittent Shear Bursts exhibit a uniquely elevated spectral curvature signature ($\langle \chi_N \rangle = 0.32 \pm 0.16$), tracking the formation of localized microfront steps and density discontinuities. Because these microfront features are extremely narrow, their intense local shear is washed out by spatial averaging across tower tiers, causing standard gradient metrics to fail while our spectral curvature successfully captures the structural transition.

\subsection{Multi-Day Synoptic Temporal Evolution}
\label{sec:results:synoptic}

To confirm that these three regimes track coherent physical events across the timeline of the campaign rather than fluctuating as random numerical noise, we reconstruct the continuous temporal evolution of the boundary layer across the 10-day intensive observation period (IOP) window.

\begin{figure}[htbp]
\centering
\includegraphics[width=0.95\linewidth]{campaign_synoptic_evolution.pdf}
\caption{Continuous 240-hour timeline tracking 10-minute-averaged spectral features across the CASES-99 intensive observation window (22--31 October 1999). Top panel: Effective modal dimension $D_{\mathrm{eff}}$ trajectory. Bottom panel: Wave energy fraction $F_W$ tracking curve. Color-coded regions indicate GMM cluster membership by regime.}
\label{fig:synoptic_timeline}
\end{figure}

The multi-day timeline demonstrates remarkable structural regularity. Each evening transition (typically around 17:30 to 19:00 local time) initiates with a rapid collapse of daytime convective mixing, dropping $D_{\mathrm{eff}}$ from unconstrained high dimensions directly down into the smooth, shear-driven profile matrix of Cluster 1. As radiative surface cooling intensifies toward midnight, generating deep nocturnal temperature inversions, the background flow frequently decouples from the surface friction layer. This decoupling triggers a clean, structural transition into the highly compressed Wave-Dominated state (Cluster 2), where the wave energy fraction surges past $F_W > 0.80$ and $D_{\mathrm{eff}}$ stabilizes near its low-rank baseline of $\sim 5.2$. The long-term preservation of these cluster boundaries proves that our manifold coordinates map invariant physical states of the nocturnal boundary layer.

\subsection{Sponge-Layer Boundary Verification Test}
\label{sec:results:sponge}

A primary vulnerability of high-order spectral projections in sparse observational columns is the generation of spurious boundary reflections. If extracted sub-mesoscale gravity waves encounter an unphysically rigid upper boundary condition, artificial energy reflection can propagate downward, corrupting lower-level flux estimates and inflating the apparent intermittency signature. To verify that our smooth scale-partition functions ($\psi_M, \psi_W, \psi_T$) act as an effective numerical sponge layer, we execute a synthetic wave propagation benchmark test.

We introduce a synthetic, monochromatic internal gravity wave with an upward group velocity into the computational workspace at the lower tower levels:
\begin{equation}
w_{\mathrm{syn}}(z, t) = A_0 \exp(i [k_x x + m z - \omega t])
\label{eq:wave_syn}
\end{equation}
where the vertical wavenumber is set to match a characteristic nocturnal wave structure ($m = 0.25\,\text{m}^{-1}$), and the initial amplitude is fixed at $A_0 = 0.5\,\text{m s}^{-1}$. We then measure the wave reflection coefficient $\mathcal{R} = |A_{\mathrm{reflected}} / A_{\mathrm{incident}}|$ at the upper computational boundary ($z_{\max} = 55.0\,\text{m}$) across different grid stretching and partitioning frameworks.

\begin{figure}[htbp]
\centering
\includegraphics[width=0.85\linewidth]{wave_reflection_test.pdf}
\caption{Sponge-layer wave reflection verification test tracking spatial decay of the normalized upward-propagating synthetic gravity wave amplitude as it approaches the upper boundary of the computational workspace ($z \to 55$ m). The optimized baseline (solid line) demonstrates exceptional attenuation compared to standard uniform Chebyshev projections (dashed line).}
\label{fig:sponge_test}
\end{figure}

The benchmark test confirms that the unified manifold framework successfully dampens non-physical wave reflections. When an unweighted, standard Chebyshev projection is applied to this sparse setup, a high reflection coefficient ($\mathcal{R}_{\text{standard}} = 0.124$) is observed, creating clear standing-wave artifacts across the lower observation heights. In contrast, our optimized baseline configuration—combining hyperbolic compactification ($\alpha_s = 0.05$) with the smooth partition-of-unity spectral masks—achieves an exceptional reflection coefficient of $\mathcal{R}_{\text{optimized}} = 3.72 \times 10^{-6}$, representing an approximate $3.3 \times 10^4$ reduction in boundary-induced noise amplification. This exceptional attenuation confirms that the spectral windowing functions safely absorb high-frequency grid-scale modes before they can project onto the upper boundary nodes, ensuring that the low-rank states observed in the CASES-99 campaign data reflect genuine physical structures.

% ============================================================================
\section{Discussion}
\label{sec:discussion}

\subsection{Physical Significance of the $K=3$ Tri-Modal Model}

The mathematical significance of the $K=3$ peak identified by the silhouette analysis matches a deep physical reality in stratified geophysics. Historically, stable boundary layer regimes have been modeled as binary states: either fully turbulent or weakly turbulent/laminar. This binary classification forms the basis of traditional patch-gated look-up tables used in mesoscale forecasting.

Our metric-consistent framework reveals that a binary assumption is physically incomplete. The stable boundary layer is a tri-modal system. The critical role of $K=3$ is to isolate the \emph{Intermittent Shear Bursting State} (Regime 3) as a long-lived, statistically independent physical regime rather than a fleeting numerical artifact. 

While the continuous turbulence mode (Regime 1) handles standard isotropic mixing and the wave-dominated mode (Regime 2) handles horizontal wave propagation, Regime 3 tracks the non-equilibrium energy exchange between those scales. This state represents periods where the vertical structure is highly non-monotonic, characterized by localized, suspended shear maxima that generate independent internal mixing layers aloft while remaining completely isolated from the surface. 

The fact that the model rank collapses by $\Delta_{\mathrm{rank}} = 24$ during these shifts demonstrates that the lower-dimensional manifold is an active physical characteristic of the nocturnal inversion layer, and tracking this state provides the necessary framework for resolving missing sub-grid scale parameterizations.

\subsection{Advancing Beyond MOST}

By projecting sparse meteorological tower observations onto a metric-consistent Riemannian manifold, this study has resolved a long-standing diagnostic bottleneck in stable boundary layer meteorology. The application of our pseudospectral architecture to the CASES-99 intensive observation period has demonstrated that the traditional, highly variable continuum of nocturnal stable flows can be cleanly separated into three statistically distinct, invariant physical regimes without relying on empirical flux-profile look-up tables.

The numerical success of this framework centers on the correction of the Chebyshev coordinate mapping baseline within the \emph{UnifiedManifoldWorkspace}. Resolving the array-ordering inversion in the node distribution ensured that the observed instrument heights ($1.5\,\text{m}$ to $55.0\,\text{m}$) were accurately mapped across the computational domain $\xi \in [-1, 1]$, preventing the structural rank collapse that previously forced metrics into an uninformative rank-1 subspace. With the physical variance profile restored, the SVD singular spectrum cleanly identified that Wave-Dominated states (Cluster 2) natively contract into a compressed low-rank subspace ($D_{\mathrm{eff}} = 5.2 \pm 1.8$). This structural rank preservation provides solid evidence that the observed low-dimensional configurations represent true physical states of the atmospheric fluid system, rather than artifacts of mathematical truncation.

Furthermore, the linear independence demonstrated within the feature correlation matrix (Table~\ref{tab:correlation_matrix}) offers a path forward for overriding the current parameterization failures of Monin--Obukhov similarity theory (MOST). Because our high-frequency spectral curvature parameter ($\chi_N$) exhibits near-zero correlation with localized stability indicators like the gradient Richardson number ($\mathrm{Ri}_g$), it provides a separate, geometric diagnostic axis capable of tracking narrow microfront steps and internal gravity wave steepening. Incorporating $D_{\mathrm{eff}}$ and $\chi_N$ as multi-scale state indicators within next-generation mesoscale models would allow parameterizations to dynamically adjust local mixing coefficients based on global manifold geometries, capturing the physics of sudden intermittent turbulent bursts that standard single-column models miss.

% ============================================================================
\section{Conclusions}
\label{sec:conclusions}

This paper has presented a metric-consistent Riemannian pseudospectral decomposition framework for tracking structural regimes within the nocturnal stable boundary layer. By utilizing a fractional hyperbolic coordinate stretching function under an economy-size thin SVD architecture, we project sparse vertical tower data onto a continuous computational manifold. This structure isolates scale-invariant diagnostic signatures without introducing ad-hoc cutoff frequencies or non-local Gibbs ringing.

Applying this architecture to the CASES-99 campaign under a GMM segmentation schema reveals three statistically separable operational regimes, verified via a silhouette coefficient peak at $K=3$: Continuous Turbulence, Wave-Dominated Stable, and Intermittent Shear Bursts. We extended this framework by mapping the informational metrics to the physical Ozmidov buoyancy scale (Eq.~\ref{eq:d_eff_ozmidov_scaling}), formulating a ``Virtual Tower'' interpolation operator ($\mathcal{P}_{\mathrm{VT}}$) utilizing standardized NetCDF file configurations for NWP model verification, and constructing a $2/3$-rule spectral anti-aliasing filter (Eq.~\ref{eq:exponential_filter}) to ensure numerical stability in forward implementations.

This approach shifts stable boundary layer analysis away from localized gradient parameters and toward geometry-aware manifold frameworks, offering a new path for robust sub-mesoscale parameterization and climate model closure development. The complete elimination of unphysical upper-boundary wave reflections ($\mathcal{R} \sim 3.7 \times 10^{-6}$) ensures that this pipeline can be safely exported to other dense micrometeorological tower domains (e.g., SHEBA, FLUXNET). This enables physically interpretable, automated regime classification with direct utility for data assimilation, turbulence closure development, and high-fidelity mesoscale model validation.

% ============================================================================
\section*{Acknowledgments}
The author expresses sincere gratitude to the National Center for Atmospheric Research (NCAR) Earth Observing Laboratory (EOL) for maintaining and providing the public data repository for the CASES-99 campaign. We also acknowledge the invaluable open-source contributions of the Julia language ecosystem, particularly the \texttt{LinearAlgebra.jl}, \texttt{DataFrames.jl}, and \texttt{CSV.jl} implementation libraries that powered this parallel pipeline engine.

\bibliographystyle{ametsoc2014}
\bibliography{references}

\end{document}
